\chapter{Unit 1: Solid State Physics}

\section{Section 1: Chapter-Wise Multiple-Choice Questions (60 MCQs)}

\subsection*{Part I: Foundational Level (Questions 1 to 20)}
\mcqitem{1}{In the free electron theory of metals, conduction electrons are treated as:}{Positive ions moving through lattice}{Free particles moving through the metal}{Waves confined inside nucleus}{Particles fixed to individual atoms}
\mcqitem{2}{Which particles mainly carry electric current in a metal according to free electron theory?}{Bound protons}{Lattice neutrons}{Positive nuclei}{Conduction electrons}
\mcqitem{3}{Drift current in a semiconductor is produced mainly by:}{A constant magnetic field}{A carrier concentration gradient}{A uniform temperature alone}{An applied electric field}
\mcqitem{4}{Diffusion current occurs when charge carriers move from a region of:}{Low concentration to high concentration}{High concentration to low concentration}{Low potential to high potential only}{Low temperature to high temperature only}
\mcqitem{5}{At absolute zero ($0\,\text{K}$), the Fermi energy is the energy of the:}{Highest empty electron state}{Lowest forbidden energy state}{Lowest occupied electron state}{Highest occupied electron state}
\mcqitem{6}{The Fermi-Dirac distribution function gives the probability that an energy state is:}{Converted into lattice energy}{Occupied by a proton}{Located in a forbidden band}{Occupied by an electron}
\mcqitem{7}{At the Fermi energy $E = E_F$, the Fermi-Dirac occupation probability is:}{$1$}{$0$}{$1/4$}{$1/2$}
\mcqitem{8}{A range of electron energies that cannot exist in a solid is called a:}{Forbidden energy band}{Valence energy band}{Permitted energy level}{Conduction energy band}
\mcqitem{9}{The highest energy band normally occupied by electrons at $0\,\text{K}$ in a semiconductor is the:}{Conduction band}{Valence band}{Vacuum band}{Forbidden band}
\mcqitem{10}{A hole in a semiconductor behaves like a particle with:}{Negative charge}{Zero charge}{Variable charge}{Positive charge}
\mcqitem{11}{Effective mass is used to describe how a charge carrier responds to:}{Nuclear force inside an atom}{An external force in a crystal}{Pressure in a vacuum}{Gravity outside the material}
\mcqitem{12}{The Hall effect is the development of a transverse voltage when a current-carrying material is placed in a:}{Longitudinal thermal field}{Uniform gravitational field}{Parallel electric field}{Perpendicular magnetic field}
\mcqitem{13}{For a material with one type of carrier having concentration $n$ and charge $q$, the Hall coefficient is:}{$R_H = q/n$}{$R_H = 1/(nq)$}{$R_H = n/q$}{$R_H = nq$}
\mcqitem{14}{A pure semiconductor without intentionally added impurities is called:}{A metallic conductor}{An extrinsic semiconductor}{An ionic insulator}{An intrinsic semiconductor}
\mcqitem{15}{Adding a pentavalent impurity to silicon usually produces:}{An intrinsic semiconductor}{An n-type semiconductor}{A perfect insulator}{An p-type semiconductor}
\mcqitem{16}{In an intrinsic semiconductor, the Fermi level lies approximately:}{At the bottom of the valence band}{Near the middle of the band gap}{Far above the conduction band}{At the top of the conduction band}
\mcqitem{17}{In an n-type semiconductor, the Fermi level shifts closer to the:}{Conduction band}{Bottom of the forbidden band}{Middle of the nucleus}{Valence band}
\mcqitem{18}{In a direct band gap semiconductor, an electron can recombine with a hole without requiring a change in:}{Electric charge}{Carrier number}{Crystal momentum}{Electron energy}
\mcqitem{19}{Which material is commonly classified as an indirect band gap semiconductor?}{Gallium nitride}{Gallium arsenide}{Silicon}{Indium phosphide}
\mcqitem{20}{A solar cell converts light energy directly into electrical energy through the:}{Piezoelectric effect}{Photovoltaic effect}{Thermoelectric effect}{Hall effect}

\subsection*{Part II: Intermediate Analytical Level (Questions 21 to 40)}
\mcqitem{21}{In the Drude model, $\sigma = \frac{ne^2\tau}{m}$. If electron density doubles while mean collision time halves, conductivity:}{Becomes half as large}{Becomes twice as large}{Remains unchanged}{Becomes four times as large}
\mcqitem{22}{Electron concentration increases along $+x$. Which electric-field direction balances electron diffusion current?}{Along the positive $x$-direction}{No electric field is required}{Perpendicular to $x$-direction}{Along the negative $x$-direction}
\mcqitem{23}{For 3D free electron gas at $0\,\text{K}$, $E_F \propto n^{2/3}$. If density increases 8 times, new Fermi energy is:}{$4 E_F$}{$8 E_F$}{$16 E_F$}{$2 E_F$}
\mcqitem{24}{Probability that an electronic state at Fermi energy $E_F$ is occupied at nonzero temperature is:}{$1$}{$0$}{$1/e$}{$1/2$}
\mcqitem{25}{At temperature $T$, a state has $E - E_F = k_B T \ln 3$. What is its Fermi-Dirac occupation probability?}{$1/2$}{$3/4$}{$1/3$}{$1/4$}
\mcqitem{26}{When $N$ identical isolated atoms form a crystal, one atomic energy level splits into:}{Exactly two levels}{About $N$ closely spaced levels}{A single level}{A forbidden gap}
\mcqitem{27}{Why does a completely filled energy band produce no net electrical current under weak field?}{Contains only positive carriers}{Electrons have no kinetic energy}{All enter conduction band}{Contributions from occupied states cancel across the band}
\mcqitem{28}{Near the top of valence band, $d^2E/dk^2 < 0$. What is the standard description of transport?}{Negative holes, negative mass}{Neutral carriers, zero mass}{Free electrons, infinite mass}{Positive holes with positive effective mass}
\mcqitem{29}{An electron occupies a band region where $m^* < 0$. If an external force acts along $+x$, its acceleration is:}{Always zero}{Along $+x$}{Along $-x$}{Perpendicular to $+x$}
\mcqitem{30}{Current is along $+x$, magnetic field along $+z$. If electrons are majority carriers, where do they accumulate?}{On the negative $y$ side}{On the positive $y$ side}{On the negative $z$ side}{On the positive $x$ side}
\mcqitem{31}{A semiconductor slab has $n = 5.0 \times 10^{22}\,\text{m}^{-3}, t = 0.50\,\text{mm}, I = 2.0\,\text{mA}, B = 0.50\,\text{T}$. Magnitude of Hall voltage $V_H$ is:}{$0.50\,\text{mV}$}{$0.25\,\text{mV}$}{$2.50\,\text{mV}$}{$1.00\,\text{mV}$}
\mcqitem{32}{Hall coefficient is $R_H = -3.9 \times 10^{-4}\,\text{m}^3/\text{C}$. The carrier type and approximate concentration are:}{Holes, $1.6 \times 10^{22}\,\text{m}^{-3}$}{Electrons, $3.9 \times 10^{22}\,\text{m}^{-3}$}{Holes, $3.9 \times 10^{22}\,\text{m}^{-3}$}{Electrons, $1.6 \times 10^{22}\,\text{m}^{-3}$}
\mcqitem{33}{In intrinsic semiconductor $n=p=n_i$. If $\mu_e = 3\mu_h$, what fraction of conductivity is due to electrons?}{$1/2$}{$1/4$}{$3/4$}{$2/3$}
\mcqitem{34}{Silicon is doped with phosphorus. After donor ionization, which carriers and fixed ions are produced?}{Free electrons and positive donor ions}{Free holes and negative donor ions}{Free electrons and negative donor ions}{Free holes and positive donor ions}
\mcqitem{35}{Why does conductivity of intrinsic semiconductor increase strongly with temperature?}{More electrons cross bandgap creating pairs}{Forbidden gap fills completely}{Valence electrons per atom increases}{Crystal becomes metal}
\mcqitem{36}{For intrinsic semiconductor, $E_i = \frac{E_C+E_V}{2} + \frac{k_BT}{2}\ln(N_V/N_C)$. If $N_C > N_V$, $E_i$ is:}{Slightly below the middle of the band gap}{Slightly above the middle of the band gap}{At conduction-band edge}{At valence-band edge}
\mcqitem{37}{At $300\,\text{K}$, electron concentration in n-type Si increases by $100\times$. How far does $E_F$ shift towards $E_C$? ($k_BT = 0.0259\,\text{eV}$)}{$0.119\,\text{eV}$}{$0.059\,\text{eV}$}{$0.460\,\text{eV}$}{$0.259\,\text{eV}$}
\mcqitem{38}{Why is direct band-gap semiconductor more efficient for light emission than indirect gap?}{Recombination requires 2 phonons}{Electron-hole recombination conserves momentum without a phonon}{Valence band has larger resistance}{No available conduction states}
\mcqitem{39}{In illuminated solar cell under short circuit, primary role of depletion field is to:}{Make carriers move to same contact}{Prevent light entering}{Separate photogenerated electrons and holes}{Create sub-bandgap photons}
\mcqitem{40}{Solar cell receives $100\,\text{mW}$ of $2.0\,\text{eV}$ light. If $80\%$ photons create collected electrons, photocurrent is:}{$40\,\text{mA}$}{$80\,\text{mA}$}{$64\,\text{mA}$}{$20\,\text{mA}$}

\subsection*{Part III: Advanced Mastery Level (Questions 41 to 60)}
\mcqitem{41}{Metal B has $2\times$ electron density and $3\times$ effective mass of metal A. If $\sigma_A = \sigma_B$, relationship between collision times is:}{$\tau_B = 6\tau_A$}{$\tau_B = \frac{1}{6}\tau_A$}{$\tau_B = \frac{3}{2}\tau_A$}{$\tau_B = \frac{2}{3}\tau_A$}
\mcqitem{42}{In nondegenerate semiconductor, $n(x) = n_0 e^{x/L}$. Which electric field makes total electron current density zero?}{$E_x = \frac{kT}{qL}$}{$E_x = \frac{qL}{kT}$}{$E_x = -\frac{kT}{qL}$}{$E_x = -\frac{qL}{kT}$}
\mcqitem{43}{Excess minority electrons in p-type Si: $\Delta n(x) = \Delta n_0 e^{-x/L_n}$ for $x>0$. With no $E$-field, directions of electron motion and diffusion current are:}{Electrons move $-x$, current $+x$}{Electrons and current both $-x$}{Electrons and current both $+x$}{Electrons move $+x$, current $-x$}
\mcqitem{44}{Two 3D free electron systems at $T=0$: System B has $n_B = 8n_A$ and $m_B^* = 2m_A^*$. What is $E_{F,B}/E_{F,A}$?}{$2$}{$8$}{$1/2$}{$4$}
\mcqitem{45}{At temp $T$, state 1 is at $E_F + k_BT\ln 3$ and state 2 is at $E_F - k_BT\ln 3$. Probabilities that state 1 is occupied and state 2 is empty are:}{Both $3/4$}{$3/4$ and $1/4$}{Both $1/4$}{$1/4$ and $3/4$}
\mcqitem{46}{A narrow group of states has constant DOS symmetric about chemical potential $\mu$. Integration limits $\gg k_BT$. Thermal excitation satisfies:}{Electrons above $\mu >$ holes below $\mu$}{Holes below $\mu >$ electrons above $\mu$}{Equality only at $T=0$}{Number of electrons excited above $\mu$ equals number of holes below $\mu$}
\mcqitem{47}{Crystal has $N$ atoms, each contributing 1 electron from nondegenerate orbital. Assuming spin degeneracy, band theory predicts:}{Contains $N/2$ states, half filled}{Contains $2N$ states, half filled}{Contains $2N$ states, completely filled}{Contains $N$ states, completely filled}
\mcqitem{48}{In nearly-free electron model, why does energy gap open at Brillouin-zone boundary?}{Pauli principle shifts states}{Periodic potential couples degenerate forward and backward waves splitting standing wave energies}{Lattice confines electron permanently}{Collisions remove states}
\mcqitem{49}{2D band dispersion: $E(k_x, k_y) = E_0 - \alpha k_x^2 + \beta k_y^2$ with $\alpha, \beta > 0$. Signs of electron effective-mass components are:}{$m_x^* > 0, m_y^* < 0$}{$m_x^* < 0, m_y^* > 0$}{$m_x^* > 0, m_y^* > 0$}{$m_x^* < 0, m_y^* < 0$}
\mcqitem{50}{Nearly full valence band has negative electron curvature near max. Why is it represented by holes with positive effective mass?}{Hole carries negative charge}{Removing electron reverses both charge and current contribution}{Removing electron changes curvature}{Hole is a proton}
\mcqitem{51}{Rectangular n-type sample carries current $I$ along $+x$ in $B$ field along $+z$, thickness $t$. Hall voltage $V_H = V(+y) - V(-y)$ is:}{$V_H = \frac{IB}{net}$}{$V_H = \frac{I}{neBt}$}{$V_H = -\frac{IBt}{ne}$}{$V_H = -\frac{IB}{net}$}
\mcqitem{52}{Weak-field Hall coefficient for electrons and holes is $R_H = \frac{p\mu_h^2 - n\mu_e^2}{q(p\mu_h + n\mu_e)^2}$. Hall voltage vanishes when:}{$p\mu_e^2 = n\mu_h^2$}{$p\mu_h = n\mu_e$}{$p\mu_h^2 = n\mu_e^2$}{$p\mu_e = n\mu_h$}
\mcqitem{53}{Compensated semiconductor has $N_D - N_A = 2 n_i$. Assuming thermal equilibrium, electron concentration $n$ is:}{$n = 2 n_i$}{$n = (1+\sqrt{2})n_i$}{$n = (\sqrt{2}-1)n_i$}{$n = \sqrt{2} n_i$}
\mcqitem{54}{Two intrinsic materials have same $E_g, T$. Material B has $N_{C,B} = 4N_{C,A}$ and $N_{V,B} = N_{V,A}/4$. How does $n_{i,B}$ compare to $n_{i,A}$?}{$n_{i,B} = n_{i,A}$}{$n_{i,B} = \frac{1}{2}n_{i,A}$}{$n_{i,B} = 4n_{i,A}$}{$n_{i,B} = 2n_{i,A}$}
\mcqitem{55}{Intrinsic semiconductor has $m_h^* = 4 m_e^*$. Relative to midgap, where is $E_i$?}{Exactly at midgap}{$\frac{1}{2}k_BT\ln 4$ above midgap}{$\frac{3}{4}k_BT\ln 4$ below midgap}{$\frac{3}{4}k_BT\ln 4$ above midgap}
\mcqitem{56}{Fully ionized n-type semiconductor ($N_D \gg n_i$). If $N_D$ increases by $100\times$, how does $E_F - E_i$ change?}{Increases by $2k_BT\ln 100$}{Remains unchanged}{Increases by $k_BT\ln 100$}{Decreases by $k_BT\ln 100$}
\mcqitem{57}{Two semiconductors have equal $E_g$, one direct, one indirect. Why direct emits light more efficiently?}{Equal numbers of states}{Conduction electrons have no effective mass}{Photons supply large momentum}{Band-edge electron-hole recombination conserves crystal momentum without a phonon}
\mcqitem{58}{Indirect gap has VB max at $k=0$ and CB min at $k=k_0 \neq 0$. Near-threshold optical absorption requires:}{Two photons absorbed}{Photon alone provides momentum $\hbar k_0$}{Photon and phonon jointly conserve energy and crystal momentum}{Impurity converts gap}
\mcqitem{59}{For ideal solar cell with ideality factor 1, $V_{oc} \approx (kT/q)\ln(I_{sc}/I_0)$. At $300\,\text{K}$, by how much does $V_{oc}$ increase when intensity rises $100\times$?}{$0.238\,\text{V}$}{$0.0259\,\text{V}$}{$0.119\,\text{V}$}{$0.0595\,\text{V}$}
\mcqitem{60}{Solar cell has $V_{oc} = 0.65\,\text{V}, I_{sc} = 40\,\text{mA}, V_m = 0.52\,\text{V}, I_m = 35\,\text{mA}$. What is its fill factor?}{$0.70$}{$0.58$}{$0.91$}{$0.80$}

\newpage
\section{Section 2: Complete Answer Key \& Technical Rationales}

\subsection*{Part A: Questions 1 to 30}
\begin{table}[htbp]
\centering
\caption{\textbf{Answer Key with Rationales (Questions 1 to 30)}}
\vspace{2mm}
\small
\begin{tabularx}{\textwidth}{l c X l c X}
\toprule
\textbf{Q\#} & \textbf{Ans} & \textbf{Technical Rationale} & \textbf{Q\#} & \textbf{Ans} & \textbf{Technical Rationale} \\
\midrule
1 & \textbf{B} & Drude theory treats valence electrons as a free classical gas. & 16 & \textbf{B} & Equal electron and hole densities position $E_i$ near midgap. \\
2 & \textbf{D} & Delocalized conduction electrons are the charge carriers in metals. & 17 & \textbf{A} & Donor electrons populate near $E_C$, shifting $E_F$ upwards. \\
3 & \textbf{D} & Drift is carrier motion driven by an applied electric field. & 18 & \textbf{C} & Extrema align at $\vec{k}=0$, conserving momentum without phonons. \\
4 & \textbf{B} & Diffusion is driven from higher to lower carrier concentration gradient. & 19 & \textbf{C} & Silicon and Germanium have indirect band gaps. \\
5 & \textbf{D} & At $0\,\text{K}$, $E_F$ is the cutoff energy of highest occupied state. & 20 & \textbf{B} & Photons create electron-hole pairs separated by p-n junction field. \\
6 & \textbf{D} & Fermi-Dirac $f(E)$ gives electronic state occupancy probability. & 21 & \textbf{C} & $\sigma \propto n \tau \implies (2)(1/2) = 1$. \\
7 & \textbf{D} & At $E = E_F$, $f(E_F) = 1/(1+e^0) = 1/2 = 50\%$ at any $T > 0\,\text{K}$. & 22 & \textbf{A} & Diffusion drives electrons toward $-x$. Electric field along $+x$ exerts force in $-x$ to create opposing drift. \\
8 & \textbf{A} & The energy gap $E_g$ contains no allowed electronic wave states. & 23 & \textbf{A} & $E_F' = (8)^{2/3} E_F = 4 E_F$. \\
9 & \textbf{B} & At $0\,\text{K}$, the valence band is completely full and conduction band is empty. & 24 & \textbf{D} & $f(E_F) = 1/(1+e^0) = 1/2$ at all $T > 0$. \\
10 & \textbf{D} & Absence of an electron behaves as an effective positive charge $+e$. & 25 & \textbf{D} & $f(E) = 1/(1 + e^{\ln 3}) = 1/(1+3) = 1/4$. \\
11 & \textbf{B} & Effective mass accounts for crystal lattice periodic potential response. & 26 & \textbf{B} & Wavefunction overlap splits $N$ degenerate states into an energy band of $N$ levels. \\
12 & \textbf{D} & Transverse Lorentz force from perpendicular $\vec{B}$ field produces $V_H$. & 27 & \textbf{D} & For every electron in state $+k$, there is an electron in $-k$, summing velocity to zero. \\
13 & \textbf{B} & By definition, $R_H = E_y / (J_x B_z) = 1/(nq)$. & 28 & \textbf{D} & Negative electron curvature is mapped to holes with $q = +e$ and $m_h^* > 0$. \\
14 & \textbf{D} & Undoped pure semiconductor material is intrinsic. & 29 & \textbf{C} & $a = F/m^*$; with $m^* < 0$, acceleration is opposite to applied force. \\
15 & \textbf{B} & Pentavalent atoms donate extra electrons, forming n-type. & 30 & \textbf{A} & $\vec{F} = -e(\vec{v}_d \times \vec{B}) = -e((-v_d\hat{x}) \times B\hat{z}) = -e(v_d B \hat{y}) = -e v_d B \hat{y}$ toward $-y$. \\
\bottomrule
\end{tabularx}
\end{table}

\subsection*{Part B: Questions 31 to 60}
\begin{table}[htbp]
\centering
\caption{\textbf{Answer Key with Rationales (Questions 31 to 60)}}
\vspace{2mm}
\small
\begin{tabularx}{\textwidth}{l c X l c X}
\toprule
\textbf{Q\#} & \textbf{Ans} & \textbf{Technical Rationale} & \textbf{Q\#} & \textbf{Ans} & \textbf{Technical Rationale} \\
\midrule
31 & \textbf{B} & $V_H = \frac{IB}{net} = \frac{(2\times 10^{-3})(0.5)}{(5\times 10^{22})(1.6\times 10^{-19})(0.5\times 10^{-3})} = 0.25\,\text{mV}$. & 46 & \textbf{D} & By particle conservation and Fermi-Dirac anti-symmetry about $\mu$. \\
32 & \textbf{D} & Negative sign $\implies$ electrons; $n = 1/(|R_H|e) = 1.6\times 10^{22}\,\text{m}^{-3}$. & 47 & \textbf{B} & With 2 spin states per orbital, band holds $2N$ electrons. $N$ electrons fill half the band (metal). \\
33 & \textbf{C} & $\sigma_e / \sigma_{total} = \mu_e / (\mu_e + \mu_h) = 3/(3+1) = 3/4$. & 48 & \textbf{B} & Bragg reflection forms standing waves with different electrostatic potential energies. \\
34 & \textbf{A} & Donating an electron leaves fixed positive donor ion $N_D^+$ and free conduction electron. & 49 & \textbf{B} & $m_x^* = \hbar^2/(-2\alpha) < 0$ and $m_y^* = \hbar^2/(+2\beta) > 0$. \\
35 & \textbf{A} & Thermal excitation across $E_g$ causes exponential carrier generation $n_i \propto e^{-E_g/2k_BT}$. & 50 & \textbf{B} & Absence of negative charge electron moving with negative acceleration behaves identically to positive charge with positive mass. \\
36 & \textbf{A} & $\ln(N_V/N_C) < 0$ when $N_C > N_V$, pulling $E_i$ below midgap. & 51 & \textbf{D} & Electrons accumulate at $-y$, so $V(-y) > V(+y)$, giving negative Hall voltage $V_H = -\frac{IB}{net}$. \\
37 & \textbf{A} & $\Delta E_F = k_BT \ln(100) = 0.0259 \times 4.605 = 0.119\,\text{eV}$. & 52 & \textbf{C} & $V_H = 0 \iff R_H = 0 \implies p\mu_h^2 = n\mu_e^2$. \\
38 & \textbf{B} & Direct transition is a 1st-order process requiring no phonon lattice interaction. & 53 & \textbf{B} & $n - n_i^2/n = 2n_i \implies n = (1+\sqrt{2})n_i$. \\
39 & \textbf{C} & Built-in electric field sweeps electrons to n-side and holes to p-side. & 54 & \textbf{A} & $N_{C,B} N_{V,B} = N_{C,A} N_{V,A} \implies n_{i,B} = n_{i,A}$. \\
40 & \textbf{A} & Photon rate $= 3.125\times 10^{17}\,\text{s}^{-1} \implies I = 0.8\times 1.6\times 10^{-19}\times 3.125\times 10^{17} = 40\,\text{mA}$. & 55 & \textbf{D} & $E_i - E_{mid} = \frac{3}{4}k_BT\ln(m_h^*/m_e^*) = \frac{3}{4}k_BT\ln 4$ above midgap. \\
41 & \textbf{C} & $\sigma = \frac{n e^2 \tau}{m} \implies \frac{2 n_A e^2 \tau_B}{3 m_A} = \frac{n_A e^2 \tau_A}{m_A} \implies \tau_B = \frac{3}{2}\tau_A$. & 56 & \textbf{C} & $E_F - E_i = k_BT\ln(N_D/n_i) \implies$ increases by $k_BT\ln 100$. \\
42 & \textbf{C} & $J_n = q n \mu_n E_x + q D_n \frac{dn}{dx} = 0 \implies E_x = -\frac{kT}{qL}$. & 57 & \textbf{D} & Direct transitions conserve $\vec{k}$ directly without 2nd-order phonon scattering. \\
43 & \textbf{D} & Electrons diffuse down gradient along $+x$. Conventional current is along $-x$. & 58 & \textbf{C} & Photon provides energy $\hbar\omega \approx E_g$; phonon supplies crystal momentum $\hbar k_0$. \\
44 & \textbf{A} & $E_F \propto \frac{n^{2/3}}{m^*} \implies \frac{(8)^{2/3}}{2} = \frac{4}{2} = 2$. & 59 & \textbf{C} & $\Delta V_{oc} = (0.0259\,\text{V})\ln(100) \approx 0.119\,\text{V}$. \\
45 & \textbf{C} & $f(E_F+\Delta) = \frac{1}{1+e^{\ln 3}} = \frac{1}{4}$. Empty prob at $E_F-\Delta$ is $1-f(E_F-\Delta) = f(E_F+\Delta) = 1/4$. & 60 & \textbf{A} & $FF = \frac{0.52 \times 35}{0.65 \times 40} = 0.70$. \\
\bottomrule
\end{tabularx}
\end{table}

\newpage
\section{Section 3: Comprehensive Theory Questions \& Worked Solutions}
\begin{theoryq}{1.1}
State the fundamental assumptions of classical Drude free electron theory and derive $\sigma = \frac{ne^2\tau}{m}$.
\end{theoryq}

\begin{theorysol}
Drude model treats valence electrons as a non-interacting classical gas colliding with fixed lattice ions with collision interval $\tau$. Under field $E$, steady state drift velocity is $v_d = -e\tau E/m$, giving current density $J = ne^2\tau E/m = \sigma E$, yielding conductivity $\sigma = \frac{ne^2\tau}{m}$.
\end{theorysol}

\vspace{3mm}

\begin{theoryq}{1.2}
Derive the mathematical expression for Hall voltage $V_H$ and explain how it determines carrier type.
\end{theoryq}

\begin{theorysol}
Under magnetic field $B\hat{z}$ and current $I\hat{x}$, Lorentz force drives carriers to $-y$ until electric Hall field $E_H$ balances magnetic force: $q E_H = q v_d B \implies E_H = v_d B$. With $v_d = \frac{I}{nqwt}$, Hall voltage is $V_H = E_H w = \frac{IB}{nqt}$. Sign of $V_H$ indicates carrier polarity.
\end{theorysol}

\vspace{3mm}

